% Work-MAT | 2. rocnik | pripravovana hodina 48
% Návaznost: uživatelské připravy.zip (hodiny 01–20), autorská učebnice na GitHubu.
% Realisticky.cz: didaktická inspirace v posloupnosti a typech úloh; zadání jsou vlastní.
\documentclass[11pt,a4paper]{article}
\input{../../sablony/pripravy/priprava.sty}
\usepackage{tikz}
\nastavPripravu{2. ročník}{Orientovaný úhel}{48}{Stupně a radiány}
\begin{document}
\blok{Motivace a návaznost}
Převádíme stupně na radiány a zpět. Nejprve změř na jednotkové kružnici čtvrt a půl otočky; až potom napiš obecný vzorec.
\poznamka{Orientační nabídka pro 45--55 minut: krátké zopakování a motivace, nový výklad, tři společné vzorové úlohy, samostatná práce. Vyber podle tempa třídy; zbývající zadání slouží jako rezerva.}
\blok{Výklad}
$180^\circ=\pi\,\mathrm{rad}$ a $360^\circ=2\pi\,\mathrm{rad}$. Převod: $\alpha_{rad}=\alpha_{deg}\frac{\pi}{180^\circ}$.
\blok{Společné příklady}
\begin{prikladReseniBox}[title={Příklad 1}]
Vyřeš, resp. urči: $150^\circ\text{ na radiány}$.
\tcblower
\[\begin{aligned}150^\circ\cdot\frac{\pi}{180^\circ}=\frac{5\pi}{6}\end{aligned}\]
\end{prikladReseniBox}
\begin{prikladReseniBox}[title={Příklad 2}]
Vyřeš, resp. urči: $\frac{5\pi}{4}\text{ na stupně}$.
\tcblower
\[\begin{aligned}\frac{5\pi}{4}\cdot\frac{180^\circ}{\pi}=225^\circ\end{aligned}\]
\end{prikladReseniBox}
\poznamka{Nepřipouštěj, aby žák psal $\pi=180$ bez jednotek. Návaznost na periodičnost bude v dalších hodinách.}
\blok{Samostatné procvičování}
\begin{zadaniReseniBox}[title={Samostatná práce 1}]
\begin{enumerate}[label=\alph*)]
\item $-45^\circ\text{ jako kladný úhel}$.
\item $390^\circ\text{ po jedné otočce}$.
\item $-120^\circ\text{ jako kladný úhel}$.
\item $765^\circ\text{ po otočkách}$.
\item $-720^\circ\text{ po otočkách}$.
\item $\frac{\pi}{2}\text{ na stupně}$.
\end{enumerate}
\tcblower
\begin{enumerate}[label=\alph*)]
\item $315°$.
\item $30°$.
\item $240°$.
\item $45°$.
\item $0°$.
\item $90°$.
\end{enumerate}
\end{zadaniReseniBox}
\blok{Dobrovolné rozšíření}
\begin{zadaniReseniBox}[title={Rozšíření 1}]
\begin{enumerate}[label=\alph*)]
\item Navrhni obdobnou úlohu s jiným úhlem a zkontroluj výsledek graficky..
\end{enumerate}
\tcblower
\begin{enumerate}[label=\alph*)]
\item $Srovnej s jednotkovou kružnicí nebo grafem.$.
\end{enumerate}
\end{zadaniReseniBox}
\blok{Úkol na doma}
Každou hodnotu zdůvodni náčrtkem jednotkové kružnice nebo orientovaným úhlem; u přesných hodnot neuváděj jen desetinné aproximace.
\begin{zadaniReseniBox}[title={Domácí úkol}]
\begin{enumerate}[label=\alph*)]
\item $\frac{5\pi}{6}\text{ na stupně}$.
\item $\frac{7\pi}{4}\text{ na stupně}$.
\item $45^\circ\text{ na radiány}$.
\item $120^\circ\text{ na radiány}$.
\item $270^\circ\text{ na radiány}$.
\item $330^\circ\text{ na radiány}$.
\item $\frac{11\pi}{6}\text{ na stupně}$.
\item $-90^\circ\text{ na radiány}$.
\end{enumerate}
\tcblower
\begin{enumerate}[label=\alph*)]
\item $150°$.
\item $315°$.
\item $\frac{\pi}{4}$.
\item $\frac{2\pi}{3}$.
\item $\frac{3\pi}{2}$.
\item $\frac{11\pi}{6}$.
\item $330°$.
\item $-\frac{\pi}{2}$.
\end{enumerate}
\end{zadaniReseniBox}
\poznamka{Domácí úkol obsahuje 8 položek; předpokládaná domácí práce přibližně 15--30 minut. Po větších blocích jsou vědomě promíchané dříve probrané typy.}
\end{document}
